% !TEX root = ../../thesis.tex

\section{Integrality}\label{sec:integrality}

Integrality of the limit varifold is an essential input to the regularity argument: it ensures that tangent cones at singular points are integral cones, which constrains them to the junction configurations (wedges, folds, and higher-order junctions) needed to derive a contradiction when the $\alpha$-structural hypothesis (Definition~\ref{defn:alpha-structural}) fails (see Section~\ref{subsec:junction-resolution}). The main result of this section is:

\begin{theorem}[Integrality of the limit varifold]\label{thm:integrality}
Let $(M^{n+1}, g)$ be a closed Riemannian manifold with $2 \leq n \leq 6$, and let $\Phi$ be an optimal nested \textup{(}ONVP\textup{)} sweepout with width $W$. Let $x_i \to x_0 \in \mathfrak{m}(\Phi)$ be a min-max sequence and $V$ the varifold limit of $|\partial^*\Omega(x_i)|$. Then $V$ is an integral $n$-varifold.
\end{theorem}

The proof splits into two cases according to whether mass cancellation occurs (Definition~\ref{def:mass-cancellation}). The no-mass-cancellation case uses the de~Lellis--Tasnady criterion (Proposition~\ref{prop:p2-DLT-integrality}); the mass-cancellation case exploits nestedness of the sweepout to count sheets directly via a graphical decomposition lemma.

\subsection{The no-mass-cancellation case}\label{subsec:integrality-no-cancel}

\begin{theorem}[Integrality without mass cancellation]\label{prop:limit-dichotomy}\label{thm:integrality-mult-one}
Assume no mass cancellation occurs at $x_0$, i.e., $\mathbf{M}(\Phi(x_0)) = W$. Then $V = |\partial^*\Omega(x_0)|$; in particular, $V$ is an integral varifold.
\end{theorem}

\begin{proof}
Since $\Phi$ is continuous in the flat topology, we have $\Omega(x_i) \to \Omega(x_0)$ in $L^1(M)$, so $D\chi_{\Omega(x_i)} \to D\chi_{\Omega(x_0)}$ weakly as measures. Moreover, the no-mass-cancellation hypothesis gives
\[
\mathrm{Per}(\Omega(x_i)) = \mathbf{M}(\Phi(x_i)) \to W = \mathbf{M}(\Phi(x_0)) = \mathrm{Per}(\Omega(x_0)).
\]
Both conditions of Proposition~\ref{prop:p2-DLT-integrality} are satisfied, so $V = |\partial^*\Omega(x_0)|$ is the integral varifold induced by the Caccioppoli boundary $\partial^*\Omega(x_0)$.
\end{proof}

\begin{figure}[ht]
    \centering
    \begin{tikzpicture}[scale=1.0]
      \def\R{2.2}
      \def\h{0.9}
      \pgfmathsetmacro{\a}{sqrt(\R*\R - \h*\h)}
      \def\ew{0.4}
      \pgfmathsetmacro{\angR}{atan2(\h, \a)}

      % Fill cap (Omega): front ellipse arc + sphere arc over top
      \fill[fillA, opacity=0.4]
        (-\a, \h) arc (180:360:{\a} and {\ew})
        arc (\angR:180-\angR:\R);

      % Sphere outline
      \draw[thick] (0,0) circle (\R);

      % Boundary curve — front (solid)
      \draw[very thick, bindA] (-\a, \h) arc (180:360:{\a} and {\ew});
      % Boundary curve — back (dashed)
      \draw[thick, bindA, dashed] (-\a, \h) arc (180:0:{\a} and {\ew});

      % Equator (V) — front (solid)
      \draw[very thick, bindB] (-\R, 0) arc (180:360:{\R} and {0.45});
      % Equator (V) — back (dashed)
      \draw[thick, bindB, dashed] (-\R, 0) arc (180:0:{\R} and {0.45});

      % Labels
      \node[bindA] at (0, 1.75) {$\Omega(x_i)$};
      \node[bindA, right] at (\a+0.1, \h-0.15) {$|\partial^*\!\Omega(x_i)|$};
      \node[bindB, right] at (\R+0.1, -0.15) {$V = |\partial^*\!\Omega(x_0)|$};
      \node[above left] at ({-\R*cos(40)}, {\R*sin(40)}) {$M$};
    \end{tikzpicture}
    \caption{The no-mass-cancellation case on a closed manifold $M$. The shaded region is the Caccioppoli set $\Omega(x_i)$, with reduced boundary $|\partial^*\!\Omega(x_i)|$ (red). As $x_i \to x_0$, the varifolds $|\partial^*\!\Omega(x_i)|$ converge to $V = |\partial^*\!\Omega(x_0)|$ (blue), which is integral with multiplicity~$1$.}
    \label{fig:caccioppoli-set}
\end{figure}

The identification $V = |\partial^*\Omega(x_0)|$ has two important consequences beyond integrality: it tells us that $V$ is the boundary of a Caccioppoli set, which constrains the possible tangent cones (Section~\ref{subsec:tangent-cones}), and it connects $V$ to the non-excessive structure of the sweepout, which provides the homotopic minimizing property.

\subsection{The mass cancellation case}\label{subsec:integrality-cancel}

When mass cancellation occurs ($\mathbf{M}(\Phi(x_0)) < W$), we still have $\Omega(x_i) \to \Omega(x_0)$ in $L^1$, but the perimeter convergence condition in Proposition~\ref{prop:p2-DLT-integrality} fails:
\[
    \mathrm{Per}(\Omega(x_i)) \to W > \mathrm{Per}(\Omega(x_0)).
\]
The mass defect $W - \mathrm{Per}(\Omega(x_0)) > 0$ is carried by sheets of $\partial^*\Omega(x_i)$ that converge in pairs, cancelling in the flat topology but contributing positively to varifold mass. In particular, $V \neq |\partial^*\Omega(x_0)|$, and the integrality of $V$ does not follow from Proposition~\ref{prop:p2-DLT-integrality}. Instead, we exploit the nestedness of the sweepout to count sheets and deduce integrality directly. The main tool is the following graphical decomposition lemma for Caccioppoli boundaries converging to a stationary varifold.

\begin{figure}[ht]
    \centering
    \begin{tikzpicture}[scale=0.75]
      \def\Ro{2.2}   % outer radius
      \def\Ri{1.2}   % inner radius

      % ==================== Left copy ====================
      \begin{scope}[xshift=-3.5cm]
        \coordinate (OT) at ({2.2*cos(50)}, {2.2*sin(50)});
        \coordinate (OB) at ({2.2*cos(310)}, {2.2*sin(310)});
        \coordinate (IT) at ({1.2*cos(50)}, {1.2*sin(50)});
        \coordinate (IB) at ({1.2*cos(310)}, {1.2*sin(310)});

        % Region 1 (upper-left): angle 110° to 160°
        \fill[fillA, opacity=0.4]
          ({2.2*cos(110)}, {2.2*sin(110)}) arc (110:160:\Ro)
          .. controls ({2.2*cos(160)+0.24*sin(160)}, {2.2*sin(160)-0.24*cos(160)})
          and ({1.2*cos(160)+0.24*sin(160)}, {1.2*sin(160)-0.24*cos(160)})
          .. ({1.2*cos(160)}, {1.2*sin(160)})
          arc (160:110:\Ri)
          .. controls ({1.2*cos(110)+0.24*sin(110)}, {1.2*sin(110)-0.24*cos(110)})
          and ({2.2*cos(110)+0.24*sin(110)}, {2.2*sin(110)-0.24*cos(110)})
          .. cycle;
        \begin{scope}[shift={({1.7*cos(110)}, {1.7*sin(110)})}, rotate=110]
          \draw[very thick, bindA] (0.5, 0) arc (0:-180:0.5 and 0.18);
          \draw[thick, bindA, dashed] (0.5, 0) arc (0:180:0.5 and 0.18);
        \end{scope}
        \begin{scope}[shift={({1.7*cos(160)}, {1.7*sin(160)})}, rotate=160]
          \draw[very thick, bindA] (0.5, 0) arc (0:-180:0.5 and 0.18);
          \draw[thick, bindA, dashed] (0.5, 0) arc (0:180:0.5 and 0.18);
        \end{scope}

        % Region 2 (lower-left): angle 200° to 250°
        \fill[fillA, opacity=0.4]
          ({2.2*cos(200)}, {2.2*sin(200)}) arc (200:250:\Ro)
          .. controls ({2.2*cos(250)+0.24*sin(250)}, {2.2*sin(250)-0.24*cos(250)})
          and ({1.2*cos(250)+0.24*sin(250)}, {1.2*sin(250)-0.24*cos(250)})
          .. ({1.2*cos(250)}, {1.2*sin(250)})
          arc (250:200:\Ri)
          .. controls ({1.2*cos(200)+0.24*sin(200)}, {1.2*sin(200)-0.24*cos(200)})
          and ({2.2*cos(200)+0.24*sin(200)}, {2.2*sin(200)-0.24*cos(200)})
          .. cycle;
        \begin{scope}[shift={({1.7*cos(200)}, {1.7*sin(200)})}, rotate=200]
          \draw[very thick, bindA] (0.5, 0) arc (0:-180:0.5 and 0.18);
          \draw[thick, bindA, dashed] (0.5, 0) arc (0:180:0.5 and 0.18);
        \end{scope}
        \begin{scope}[shift={({1.7*cos(250)}, {1.7*sin(250)})}, rotate=250]
          \draw[very thick, bindA] (0.5, 0) arc (0:-180:0.5 and 0.18);
          \draw[thick, bindA, dashed] (0.5, 0) arc (0:180:0.5 and 0.18);
        \end{scope}

        % Outline
        \draw[thick] (OT) arc (50:310:\Ro);
        \draw[thick] (IB) arc (310:50:\Ri);
        \draw[thick] (OB) to[out=40, in=40, looseness=1.5] (IB);
        \draw[thick] (IT) to[out=320, in=320, looseness=1.5] (OT);

        \node at (0.5, 0) {$M$};
      \end{scope}

      % ==================== Right copy ====================
      \begin{scope}[xshift=3.5cm]
        \coordinate (OT2) at ({2.2*cos(50)}, {2.2*sin(50)});
        \coordinate (OB2) at ({2.2*cos(310)}, {2.2*sin(310)});
        \coordinate (IT2) at ({1.2*cos(50)}, {1.2*sin(50)});
        \coordinate (IB2) at ({1.2*cos(310)}, {1.2*sin(310)});

        % Single merged region: angle 110° to 250°
        \fill[fillA, opacity=0.4]
          ({2.2*cos(110)}, {2.2*sin(110)}) arc (110:250:\Ro)
          .. controls ({2.2*cos(250)+0.24*sin(250)}, {2.2*sin(250)-0.24*cos(250)})
          and ({1.2*cos(250)+0.24*sin(250)}, {1.2*sin(250)-0.24*cos(250)})
          .. ({1.2*cos(250)}, {1.2*sin(250)})
          arc (250:110:\Ri)
          .. controls ({1.2*cos(110)+0.24*sin(110)}, {1.2*sin(110)-0.24*cos(110)})
          and ({2.2*cos(110)+0.24*sin(110)}, {2.2*sin(110)-0.24*cos(110)})
          .. cycle;
        \begin{scope}[shift={({1.7*cos(110)}, {1.7*sin(110)})}, rotate=110]
          \draw[very thick, bindA] (0.5, 0) arc (0:-180:0.5 and 0.18);
          \draw[thick, bindA, dashed] (0.5, 0) arc (0:180:0.5 and 0.18);
        \end{scope}
        \begin{scope}[shift={({1.7*cos(250)}, {1.7*sin(250)})}, rotate=250]
          \draw[very thick, bindA] (0.5, 0) arc (0:-180:0.5 and 0.18);
          \draw[thick, bindA, dashed] (0.5, 0) arc (0:180:0.5 and 0.18);
        \end{scope}
        % Blue ellipse at center (180°) — mass cancellation locus
        \begin{scope}[shift={({1.7*cos(180)}, {1.7*sin(180)})}, rotate=180]
          \draw[very thick, bindB] (0.5, 0) arc (0:-180:0.5 and 0.18);
          \draw[thick, bindB, dashed] (0.5, 0) arc (0:180:0.5 and 0.18);
        \end{scope}

        % Outline
        \draw[thick] (OT2) arc (50:310:\Ro);
        \draw[thick] (IB2) arc (310:50:\Ri);
        \draw[thick] (OB2) to[out=40, in=40, looseness=1.5] (IB2);
        \draw[thick] (IT2) to[out=320, in=320, looseness=1.5] (OT2);

        \node at (0.5, 0) {$M$};
      \end{scope}
    \end{tikzpicture}
    \caption{A closed manifold $M \cong S^2$ bent into a C-shape. \textbf{Left:} two disjoint sweepout slices $\partial^*\Omega(x_i)$ (red) with no mass cancellation. \textbf{Right:} the slices extend and overlap; the blue cross-section marks where the two sheets meet with opposite orientations, so the mass cancels in the varifold limit.}
    \label{fig:mass-cancellation-manifold}
\end{figure}

\begin{lemma}[Sheet decomposition]\label{lem:sheet-decomposition}
Let $(M^{n+1}, g)$ be a closed Riemannian manifold with $2 \leq n \leq 6$. Let $\{\Omega_i\}$ be a sequence of Caccioppoli sets with $|\partial^*\Omega_i| \to V$ as varifolds, where $V$ is a stationary rectifiable $n$-varifold. Suppose that $p \in \spt\|V\|$ is a point where $V$ has a unique tangent plane $P = T_p\Sigma$ and finite density $\theta(p) = \Theta(\|V\|, p) \in (0, \infty)$ (these conditions hold at $\mathcal{H}^n$-a.e.\ point of $\Sigma$ by rectifiability; see Section~\ref{sec:p2-prelim} and~\cite[Section~11]{Sim84}).

Then for every $\varepsilon > 0$, there exist $r_0 > 0$ and $i_0 \in \mathbb{N}$ such that for a.e.\ $r \in (0, r_0)$ and all $i \geq i_0$, the following holds:
\begin{enumerate}[label=\textup{(\roman*)}]
    \item $\partial^*\Omega_i \cap B_r(p)$ consists of exactly $k$ smooth connected components $\Sigma_i^1, \ldots, \Sigma_i^k$, where $k = k(r)$ is independent of $i$ and satisfies $\theta(p) = k$;
    \item each $\Sigma_i^j$ is the graph of a smooth function $u_i^j \colon \Omega_i^j \to \mathbb{R}$ over a domain $\Omega_i^j \subset D_r$ with $\|Du_i^j\|_{C^0} < \varepsilon$ and $\mathcal{H}^n(D_r \setminus \Omega_i^j) < \varepsilon \omega_n r^n$, where $D_r = P \cap B_r(p)$;
    \item the area of each sheet satisfies $|\mathcal{H}^n(\Sigma_i^j) - \omega_n r^n| < \varepsilon \omega_n r^n$.
\end{enumerate}
In particular, $\theta(p) \in \mathbb{Z}_{>0}$.
\end{lemma}

\begin{figure}[ht]
    \centering
    \begin{minipage}[b]{0.58\textwidth}
      \centering
      \begin{tikzpicture}[xshift=-0.8cm]
        % Junction point
        \coordinate (J) at (0,0);
        % Long arm: slight upward curve to the right
        \draw[thick] (J) to[out=10, in=170] coordinate[pos=0.75] (P) (3, 0.4);
        % Tangent line at P (blue)
        \draw[thick, bindB] ($(P) + (-1.2, -0.12)$) node[left] {$P = T_p\Sigma$} -- ($(P) + (1.2, 0.12)$);
        % Two short arms going upper-left and lower-left
        \draw[thick] (J) to[out=140, in=0] (-1.8, 1.2);
        \draw[thick] (J) to[out=220, in=0] (-1.8, -1.2);
        % Junction point dot
        \fill (J) circle (2pt);
        % Red sheets — thin bright red (outside circle)
        \draw[thin, red] (-2.2, 1.6) to[out=-10, in=170] ($(P) + (-0.8, 0.3)$) to[out=-10, in=170] ($(P) + (0.8, 0.25)$) to[out=-10, in=170] (3.5, 0.6);
        \draw[thin, red] (-2.2, -1.6) to[out=10, in=190] ($(P) + (-0.8, -0.25)$) to[out=-5, in=175] ($(P) + (0.8, -0.25)$) to[out=-5, in=175] (3.5, -0.3);
        \draw[thin, red] (-1.5, 1.0) to[out=350, in=90] (-0.8, 0) to[out=270, in=10] (-1.5, -1.0);
        % Red sheets — thick (inside circle, clipped)
        \begin{scope}
          \clip (P) circle (0.8);
          \draw[very thick, bindA] (-2.2, 1.6) to[out=-10, in=170] ($(P) + (-0.8, 0.3)$) to[out=-10, in=170] coordinate[pos=0.5] (S1) ($(P) + (0.8, 0.25)$) to[out=-10, in=170] (3.5, 0.6);
          \draw[very thick, bindA] (-2.2, -1.6) to[out=10, in=190] ($(P) + (-0.8, -0.25)$) to[out=-5, in=175] coordinate[pos=0.5] (S2) ($(P) + (0.8, -0.25)$) to[out=-5, in=175] (3.5, -0.3);
        \end{scope}
        % Cut curves at B_r(p) boundary: white eraser then gray circle
        \draw[white, line width=4pt] (P) circle (0.8);
        \draw[gray, thick] (P) circle (0.8);
        \node[gray, above right] at ($(P) + (0.55, 0.55)$) {$B_r(p)$};
        % Blue point on the long arm (drawn on top)
        \fill[bindB] (P) circle (2.5pt);
        % Labels for the two connected components — on top layer
        \node[bindA] (L1) at ($(P) + (-0.2, 1.4)$) {$\Sigma_i^1$};
        \draw[bindA, ->] (L1) -- (S1);
        \node[bindA] (L2) at ($(P) + (-0.2, -1.4)$) {$\Sigma_i^2$};
        \draw[bindA, ->] (L2) -- (S2);
        % Label for all three red curves
        \node[red] at (-2.8, 0) {$\partial^*\Omega(x_i)$};
      \end{tikzpicture}\\[4pt]
      \textbf{(a)}
    \end{minipage}%
    \hfill
    \begin{minipage}[b]{0.38\textwidth}
      \centering
      \begin{tikzpicture}[scale=1.5]
        % Define P at origin for this zoomed-in view
        \coordinate (P) at (0,0);
        % Named paths for intersection computation
        \path[name path=circ] (P) circle (0.8);
        \path[name path=sheet1] (-1.5, 0.5) to[out=-3, in=183] (1.5, 0.44);
        \path[name path=sheet2] (-1.5, -0.62) to[out=3, in=177] (1.5, -0.68);
        % Find intersections
        \path[name intersections={of=circ and sheet1, by={I1L,I1R}}];
        \path[name intersections={of=circ and sheet2, by={I2L,I2R}}];
        % Gray circle
        \draw[gray, thick] (P) circle (0.8);
        \node[gray, above right] at ($(P) + (0.55, 0.55)$) {$B_r(p)$};
        % Blue tangent line (clipped to circle) — horizontal
        \begin{scope}
          \clip (P) circle (0.8);
          \draw[thick, bindB] ($(P) + (-1.2, 0)$) -- ($(P) + (1.2, 0)$);
        \end{scope}
        % Two thick red sheets (clipped to circle)
        \begin{scope}
          \clip (P) circle (0.8);
          \draw[very thick, bindA] (-1.5, 0.5) to[out=-3, in=183] coordinate[pos=0.5] (S1b) (1.5, 0.44);
          \draw[very thick, bindA] (-1.5, -0.62) to[out=3, in=177] coordinate[pos=0.5] (S2b) (1.5, -0.68);
        \end{scope}
        % Vertical dashed red lines from exact intersections to D_r
        \draw[thin, bindA, dashed] (I1L) -- (I1L |- P);
        \draw[thin, bindA, dashed] (I1R) -- (I1R |- P);
        \draw[thin, bindA, dashed] (I2L) -- (I2L |- P);
        \draw[thin, bindA, dashed] (I2R) -- (I2R |- P);
        % Braces for domains on D_r
        \draw[red, thick, decorate, decoration={brace, mirror, amplitude=4pt}]
          (I1L |- P) -- (I1R |- P);
        \node[red] (OL2) at (1.2, -0.25) {$\Omega_i^2$};
        \draw[red, thin, ->] (OL2) to[bend left=40] (0.15, -0.07);
        \draw[red, thick, decorate, decoration={brace, amplitude=4pt}]
          (I2R |- P) -- (I2L |- P);
        \node[red] (OL1) at (-1.2, 0.25) {$\Omega_i^1$};
        \draw[red, thin, ->] (OL1) to[bend left=40] (-0.15, 0.07);
        % Labels outside circle with arrows
        \node[bindA] (L1b) at ($(P) + (-0.2, 1.3)$) {$\mathrm{graph}(u_i^1)$};
        \draw[bindA, ->] (L1b) -- (S1b);
        \node[bindA] (L2b) at ($(P) + (-0.2, -1.3)$) {$\mathrm{graph}(u_i^2)$};
        \draw[bindA, ->] (L2b) -- (S2b);
        % Tangent plane label
        \node[bindB, right] at ($(P) + (0.85, 0)$) {$D_r$};
      \end{tikzpicture}\\[4pt]
      \textbf{(b)}
    \end{minipage}
    \caption{Sheet decomposition (Lemma~\ref{lem:sheet-decomposition}) with $k = \theta(p) = 2$. \textbf{(a)}~Global view: $\partial^*\Omega(x_i)$ (red) near a point $p$ on the support of $V$, with the ball $B_r(p)$ cutting the boundary into $k = 2$ sheets $\Sigma_i^1, \Sigma_i^2$. \textbf{(b)}~Zoomed view inside $B_r(p)$: each sheet $\Sigma_i^j$ is the graph of a smooth function $u_i^j$ over the tangent plane $P = T_p\Sigma$, with $\|Du_i^j\|_{C^0} < \varepsilon$.}
    \label{fig:sheet-decomposition}
\end{figure}

\begin{proof}
The proof combines three ingredients: the monotonicity formula (Proposition~\ref{prop:p2-monotonicity}), graphical decomposition under small tilt-excess (Definition~\ref{def:p2-tilt-excess}), and the sphere constraint to exclude partial sheets.

\medskip\noindent\textbf{Step~1: Setup.}
Working in normal coordinates centered at $p$, we identify a neighborhood of $p$ with a neighborhood of the origin in $\mathbb{R}^{n+1}$, the tangent plane $P$ with $\mathbb{R}^n \times \{0\}$, and write $D_r = \{y \in \mathbb{R}^n : |y| < r\}$ for the $n$-disk in $P$. Since $V$ has tangent varifold $\theta(p)|P|$ at $p$, the density ratio satisfies
\[
    \omega_n^{-1} r^{-n} \|V\|(B_r(p)) \to \theta(p) \quad \text{as } r \to 0.
\]

\medskip\noindent\textbf{Step~2: Tilt-excess control.}
Since $V$ has tangent plane $P$ at $p$, the tilt-excess (Definition~\ref{def:p2-tilt-excess}) satisfies $E(V, B_r(p), P) \to 0$ as $r \to 0$. The varifold convergence $|\partial^*\Omega_i| \to V$ (Definition~\ref{def:p2-varifold-convergence}) implies $E(|\partial^*\Omega_i|, B_r(p), P) \to E(V, B_r(p), P)$ for each generic $r$ (with $\|V\|(\partial B_r(p)) = 0$). Combining: for every $\delta > 0$, there exist $r_1 > 0$ and $i_1$ such that for $r < r_1$ and $i > i_1$,
\[
    E(|\partial^*\Omega_i|, B_r(p), P) < \delta.
\]

\medskip\noindent\textbf{Step~3: Graphical decomposition.}
Since each $\partial^*\Omega_i$ is a smooth embedded hypersurface (Proposition~\ref{prop:p2-federer-regularity}), we apply the following standard result (cf.~\cite[Chapter~7]{Sim84} and~\cite[Lemma~2.1]{Wic14}): there exists $\delta_0 = \delta_0(n) > 0$ such that if $\Sigma$ is a smooth closed embedded hypersurface in $B_R(0) \subset \mathbb{R}^{n+1}$ satisfying
\begin{enumerate}[label=\textup{(H\arabic*)}]
    \item $\mathcal{H}^n(\Sigma \cap B_R(0)) \leq C \omega_n R^n$ for some constant $C$;
    \item $R^{-n} \int_{\Sigma \cap B_R(0)} |\nu_\Sigma - e_{n+1}|^2 \, d\mathcal{H}^n < \delta_0$;
\end{enumerate}
then for a.e.\ $\rho \in (R/2, R)$, the intersection $\Sigma \cap B_\rho(0)$ consists of finitely many connected components, each the graph of a smooth function $u \colon \Omega_u \to \mathbb{R}$ over a domain $\Omega_u \subset P$ with $\|Du\|_{C^0} \leq C(n)\delta_0^{1/2}$. (The proof uses Chebyshev's inequality to get pointwise control of $|\nu_\Sigma - e_{n+1}|$ outside a small set, then the implicit function theorem to decompose each nearly-flat piece into a graph; embeddedness prevents the graphical pieces from merging or overlapping.)

We apply this with $\Sigma = \partial^*\Omega_i$, $R = 2r$, and $\delta_0$ chosen so that $C(n)\delta_0^{1/2} < \varepsilon$. Condition~(H1) follows from varifold convergence, and condition~(H2) from the tilt-excess bound. Relabeling $\rho$ as $r$ gives, for $r$ small and $i$ large, a decomposition of $\partial^*\Omega_i \cap B_r(p)$ into finitely many smooth graphs $\Sigma_i^1, \ldots, \Sigma_i^{k_i}$ with $u_i^j \colon \Omega_i^j \to \mathbb{R}$ and $\|Du_i^j\|_{C^0} < \varepsilon$.

\medskip\noindent\textbf{Step~4: Each sheet nearly fills $D_r$, and the sheet count is constant.}
Since $\partial^*\Omega_i$ is compact without boundary (Proposition~\ref{prop:p2-federer-regularity}), the boundary of each graphical piece $\Sigma_i^j$ lies on $\partial B_r(p)$: every $y \in \partial\Omega_i^j$ satisfies $|y|^2 + |u_i^j(y)|^2 = r^2$. Consider a sheet passing through $B_{r/2}(p)$, i.e., there exists $y_0 \in \Omega_i^j$ with $|y_0| < r/2$. Since $V$ has tangent varifold $\theta(p)|P|$ at $p$, the mass of $V$ concentrates near $P$: for any $\alpha > 0$, $\|V\|(\{q \in B_r(p) : \mathrm{dist}(q, P) > \alpha r\}) / (\omega_n r^n) \to 0$ as $r \to 0$. By varifold convergence, $|u_i^j(y_0)| \leq \varepsilon r$ for $i$ large. For any $y \in \partial\Omega_i^j$, the gradient bound gives $|u_i^j(y)| \leq 3\varepsilon r$, and the sphere constraint forces $|y| \geq r(1 - C'\varepsilon^2)$. Therefore $\mathcal{H}^n(D_r \setminus \Omega_i^j) \leq C''\varepsilon^2 \omega_n r^n$, establishing~(ii). Each contributing sheet then has area satisfying $|\mathcal{H}^n(\Sigma_i^j) - \omega_n r^n| < C'''\varepsilon^2\omega_n r^n$, establishing~(iii). (Sheets not passing through $B_{r/2}(p)$ contribute vanishingly small mass near $p$ for generic $r$.)

The total area satisfies $k_i(1 - \varepsilon')\omega_n r^n < \mathcal{H}^n(\partial^*\Omega_i \cap B_r(p)) < k_i(1 + \varepsilon')\omega_n r^n$, where $\varepsilon' = C(\varepsilon^2)$, and converges to $\theta(p)\omega_n r^n$. Choosing $\varepsilon'$ small enough (so that intervals for consecutive integers are disjoint), $k_i$ is uniquely determined for $i$ large, hence eventually constant equal to $k = \theta(p) \in \mathbb{Z}_{>0}$, establishing~(i).
\end{proof}

We now apply the sheet decomposition lemma to prove integrality.

\begin{proposition}\label{prop:integrality-nested}
Let $(M^{n+1}, g)$ be a closed Riemannian manifold with $2 \leq n \leq 6$, and let $\Phi$ be an optimal nested \textup{(}ONVP\textup{)} sweepout with width $W$. Let $x_i \to x_0 \in \mathfrak{m}(\Phi)$ be a min-max sequence, $\Omega(x_0) = \lim \Omega(x_i)$ in $L^1$, and $V$ the varifold limit of $|\partial^*\Omega(x_i)|$. Then $V$ is an integral $n$-varifold.
\end{proposition}

\begin{proof}
Write $V = |\partial^*\Omega(x_0)| + V'$, where $V'$ is the non-negative varifold capturing the cancelling mass.

\medskip\noindent\textbf{Step~1: Rectifiability and densities.}
Since $V$ is stationary (Proposition~\ref{thm:CLS-stationary}), Allard's rectifiability theorem~\cite[Theorem~5.5(1)]{All72} gives that $V = \theta \, \mathcal{H}^n \mres \Sigma$ for an $\mathcal{H}^n$-rectifiable set $\Sigma$ and a function $\theta \colon \Sigma \to (0, \infty)$. The monotonicity formula (Proposition~\ref{prop:p2-monotonicity}) ensures $\theta(p) > 0$ at every $p \in \spt\|V\|$.

\medskip\noindent\textbf{Step~2: Integer density.}
By Lemma~\ref{lem:sheet-decomposition} applied with $\Omega_i = \Omega(x_i)$, at $\mathcal{H}^n$-a.e.\ point $p$ of $\spt\|V\|$, the density $\theta(p)$ equals the sheet count $k \in \mathbb{Z}_{>0}$. In particular, $\theta(p) \in \mathbb{Z}_{>0}$ a.e. It remains to determine which positive integers arise; the parity is controlled by nestedness.

\medskip\noindent\emph{Case~1: Pure cancellation points ($p \in \spt\|V'\| \setminus \spt\|\partial^*\Omega(x_0)\|$).}
At such a point, $\Omega(x_0)$ is locally constant near $p$. Say $B_r(p) \subset \Omega(x_0)$. If $x_i > x_0$, nestedness gives $B_r(p) \subset \Omega(x_i)$ and $\partial^*\Omega(x_i) \cap B_r(p) = \varnothing$; thus only parameters $x_i < x_0$ contribute sheets. For such $x_i$, $\Omega(x_i) \subset \Omega(x_0)$, so $\Omega(x_i) \cap B_r(p) = B_r(p)$ minus thin slabs, each bounded by two sheets. Therefore $k$ is even. (The case $B_r(p) \subset M \setminus \overline{\Omega(x_0)}$ is analogous.) The geometric mechanism is \emph{folding}: each pair of sheets bounds a slab that collapses as $x_i \to x_0$ (by volume parametrization: $\mathrm{Vol}(\Omega(x_i) \triangle \Omega(x_0)) = |x_i - x_0| \to 0$), contributing multiplicity~$2$.

\medskip\noindent\emph{Case~2: Mixed regular points ($p \in \mathrm{reg}(\partial^*\Omega(x_0)) \cap \spt\|V'\|$).}
Here $\partial^*\Omega(x_0)$ is smooth near $p$ (Proposition~\ref{prop:p2-federer-regularity}). By Lemma~\ref{lem:sheet-decomposition}, the $k$ sheets are ordered by height: $u_i^1 < \cdots < u_i^k$. By a volume comparison argument---since $\Omega(x_i)$ occupies alternating slabs and $\Omega(x_i) \to \Omega(x_0)$ in $L^1$---exactly one sheet (the \emph{tracking sheet}) satisfies $\|u_i^{j_0} - u_0\|_{L^\infty} \to 0$, where $u_0$ parametrizes $\partial^*\Omega(x_0)$. The remaining $k - 1$ cancelling sheets come in pairs by the same slab argument as Case~1. Therefore $k - 1$ is even, i.e., $k$ is odd.

\medskip\noindent\emph{Case~3: Singular points of $\partial^*\Omega(x_0)$.}
For $n \leq 6$, Federer regularity (Proposition~\ref{prop:p2-federer-regularity}) gives $\partial\Omega(x_0) = \partial^*\Omega(x_0)$, so this case is empty.

\medskip\noindent\textbf{Step~3: Conclusion.}
Combining Cases~1--3, $\theta(p) \in \mathbb{Z}_{>0}$ at $\mathcal{H}^n$-a.e.\ point. Together with rectifiability, $V$ is integral.
\end{proof}

\begin{remark}\label{rem:sheet-decomposition}
The sheet decomposition lemma (Lemma~\ref{lem:sheet-decomposition}) is a general result about Caccioppoli boundaries converging to a stationary varifold, independent of the min-max context. It is closely related to the Sheeting Theorem in Wickramasekera's regularity theory~\cite{Wic14}, which provides a multi-valued graphical decomposition for stable integral varifolds. The difference is that our lemma applies to the \emph{approximating} varifolds $|\partial^*\Omega_i|$ (multiplicity-$1$, smooth for $n \leq 6$) rather than the limit $V$, avoiding the need to assume integrality a priori. The non-excessive hypothesis is not used in the proof of integrality; it enters only in the regularity step (Section~\ref{sec:regularity}).
\end{remark}
